% !PRINTED 168
\weberpage{168}

Now if $w$ is any indecomposable function which divides neither
all the $u_r$ nor all the $v_s$, then we choose $r$ and $s$ in (3)
so that $u_r$ is the first $u$ not divisible by $w$, so that
$u_{r-1}, u_{r-2}\dots$ are divisible by $w$, and likewise so that $v_s$ is the
first $v$ not divisible by $w$. Then $U_{r+s}$ also cannot
be divisible by $w$, because all its terms, with the exception
of the first, are divisible by $w$; that is, $F(t)$ is primitive.

From this it follows immediately, if $p$ and $q$ are any integral
rational functions of the $x$, that $pq$ is the divisor of the product
of the two imprimitive functions $pf(t)$, $q\varphi(t)$, and hence
the second part of Theorem V. From this it follows further:

\gesperrt{VI. If an integral rational function. $F(t)$ of the $n+1$ variables $x$ and $t$ is decomposable into two factors which are integral with respect to $t$ and at least rational with respect to the $x$, then it is also decomposable into two functions which are integral and rational in $x$ and $t$.} % [sic] printed period after "function" before F(t)

For by the hypothesis there is an integral rational
function $w$ of the $x$ alone and two integral rational functions
of the $x$ and $t$, $f_1(t)$, $\varphi_1(t)$, such that
\[
w\,F(t) = f_1(t)\,\varphi_1(t),
\]
and by (V) $w$ must divide the product of the divisors of $f_1(t)$ and
$\varphi_1(t)$, so that we also obtain
\[
F(t) = f(t)\,\varphi(t)
\]
where $f(t)$, $\varphi(t)$ are likewise integral rational functions
of $x$ and $t$, and indeed with respect to $t$ of the same degree as
$f_1(t)$ and $\varphi_1(t)$, q. e. d. --- From this we conclude further
that two functions $F(t)$, $f(t)$ which, regarded as integral rational functions
of the $n+1$ variables $x$ and $t$, are relatively prime in the sense
defined above, must also prove to be relatively prime when regarded
as functions of $t$ alone and treated by the algorithm of the greatest
common divisor (\S.\,6). For if they had a common divisor which
was integral with respect to $t$ but fractional with respect to the $x$,
then an integral function $T$ of $x$ and $t$ and an integral function
$P$ of the $x$ alone could be determined so that
\[
P\,F(t) = T\,F_1(t), \qquad P\,f(t) = T\,f_1(t)
\]
would hold, where $F_1, f_1$ are integral functions without a common divisor
%CONT>>
